\section{Motivation and Background}
\label{gen:sec:motivation}

\begin{comment}
    
\subsection{Recurring Resource-Intensive Queries} 

\genjie{ Enterprise data analytics often consists of recurring queries over evolving data, which, despite being essential to daily operations, can be highly resource-intensive. In fact, over 60\% of the jobs in Microsoft clusters have been reported to be recurrent~\cite{jindal2018selecting}. Since such queries often run on a fixed schedule (e.g., nightly or multiple times per day), suboptimal SQL logic can accumulate into significant performance bottlenecks and increased cloud spend over time. Thus, identifying, analyzing, and optimizing these heavy queries is critical to sustaining an efficient data workflow.}



\genjie{Real-world dbt case studies~\cite{dbtcasestudy1,dbtcasestudy2} illustrate how data teams manually refactored complex SQL queries—executed four times per day—by removing redundancies, streamlining joins, and leveraging incremental models—efforts requiring deep domain knowledge and focused engineering expertise. Although these refactorings often dramatically reduce query runtimes, the process is nontrivial and labor-intensive. Consequently, an automated and interpretable solution is needed to help systematically address these recurring performance bottlenecks.} 

% TODO (Phase 2): the exported paper source contained a leftover author
% note here: "Add a problem definition here for our query rewriting
% setting". A shared problem definition is planned for the dissertation's
% Background/Introduction; revisit this spot when de-duplicating.


\end{comment}

\subsection{Limitations of Traditional Query Rewriting}

Despite decades of effort by experts to craft extensive rewrite rules, many queries still cannot be optimized by existing rules~\cite{dong2023slabcity}. 
In those cases, leveraging general knowledge becomes essential in identifying redundant calculations in the query and exploring alternative ways of achieving the same output. 
To demonstrate this, we present two examples: one from LeetCode~\cite{gen_leetcode}, the other from TPC-DS~\cite{tpcds}, an industry-standard business intelligence benchmark.


\begin{comment}
    

\begin{figure}[h]
  \centering
  \begin{minipage}{.45\textwidth}
    \begin{minted}[frame=single, linenos, fontsize=\footnotesize, breaklines]{sql}
      select max(a.salary) as secondHighest
      from employee as a, employee as b
      where a.salary < b.salary
    \end{minted}
    \vspace{-4mm}
     \captionsetup{font=small,name=Figure}
    \caption*{Q1: Cartesian product}
    \vspace{3mm}
  \end{minipage}
  % \vspace{2mm}
  \hfill
  \begin{minipage}{.45\textwidth}
    \begin{minted}[frame=single, linenos, fontsize=\footnotesize, breaklines]{sql}
      select max(salary) as secondHighest
      from employee
      where salary < (select max(salary) from employee)
    \end{minted}
    \vspace{-4mm}
     \captionsetup{font=small,name=Figure}
    \caption*{Q2: with subquery}
  \end{minipage}
  \vspace{-3mm}
  \captionsetup{font=small,name=Figure}
  \caption{Finding the second highest salary}
  \inv
  \label{gen:fig:sql-comparison}
\end{figure}
\end{comment}




\begin{figure}[h]
  \centering
  % Left Column: Q1
  \begin{minipage}{.47\textwidth}
    \begin{lstlisting}[
        language=SQL, 
        frame=single, 
        numbers=left, 
        basicstyle=\footnotesize\ttfamily, 
        breaklines=true,
        xleftmargin=2em, % Adjusts for line numbers inside the frame
        framexleftmargin=1.5em
    ]
select max(a.salary) as secondHighest
from employee as a, employee as b
where a.salary < b.salary
    \end{lstlisting}
    \vspace{-2mm}
    \captionsetup{font=small}
    \caption*{Q1: Cartesian product}
  \end{minipage}
  \hfill
  % Right Column: Q2
  \begin{minipage}{.47\textwidth}
    \begin{lstlisting}[
        language=SQL, 
        frame=single, 
        numbers=left, 
        basicstyle=\footnotesize\ttfamily, 
        breaklines=true,
        xleftmargin=2em,
        framexleftmargin=1.5em
    ]
select max(salary) as secondHighest
from employee
where salary < (select max(salary) from employee)
    \end{lstlisting}
    \vspace{-2mm}
    \captionsetup{font=small}
    \caption*{Q2: with subquery}
  \end{minipage}
  
  \vspace{1mm}
  \captionsetup{font=small}
  \caption{Finding the second highest salary}
  \inv
  \label{gen:fig:sql-comparison}
\end{figure}






\genhead{Example 1: LeetCode \#176} The goal is to find the second highest salary from the Employee table. \query1 in Figure~\ref{gen:fig:sql-comparison} is the human-written solution (submitted by LeetCode users) while \query2 emerges as a rewrite generated by LLMs. \query2 is approximately 2000x faster than \query1 when executed on a database  populated with one million rows. 
The optimization process involves a deep understanding of the query intent through reasoning: (1) identifying that the goal of \query1 is to find the second-highest salary, (2) realizing that there is an alternative way to achieve the same thing, i.e., first find the maximum salary using a subquery and then compare all salaries with this maximum salary to find the second maximum, and finally (3) recognizing that \query2 achieves performance improvement by reducing the number of comparisons required. No existing rewrite rule explicitly captures this transformation. As a result, this optimization cannot be achieved using traditional rule-based query rewriting techniques, highlighting the advantage of LLM-guided rewriting in discovering non-trivial, high-impact transformations.

\begin{comment}
    


\begin{listing}[ht]
  \tiny  % Reduce font size (use \footnotesize, \scriptsize, or even \tiny)
  \centering
  \begin{minipage}[t]{0.49\linewidth}  % First half
    \begin{tcolorbox}[colback=gray!5, colframe=gray!40, boxsep=2mm, left=6pt, right=5pt, top=0.5mm, bottom=0.5mm]
      \begin{minted}[
          frame=none,
          framesep=2mm,
          linenos,
          numbersep=2mm,
          escapeinside=||,
          baselinestretch=1.0, 
          fontsize=\tiny
      ]{sql}
with year_total as (
    select ..., 's' sale_type
    from  |\textcolor{blue}{$store\_sales$}|, ...
    ...
    union all
    select ..., 'w' sale_type
    from  |\textcolor{blue}{$web\_sales$}|, ...
    ...
)
select
    s_2.customer_id, ...
from
    year_total s_1,
    year_total s_2,
    year_total w_1,
    year_total w_2
where
    s_2.customer_id = s_1.customer_id
    and s_1.customer_id = w_2.customer_id
    and s_1.customer_id = w_1.customer_id
    and s_1.sale_type = 's'
    and s_1.dyear = 1998
    and w_1.sale_type = 'w'
    and w_1.dyear = 1998
    and s_2.sale_type = 's'
    and s_2.dyear = 1999
    and w_2.sale_type = 'w'
    and w_2.dyear = 1999
    and w_2.y_total / w_1.y_total > 
        s_2.y_total / s_1.y_total
      \end{minted}
    \end{tcolorbox}
    \vspace{-1mm}
    \caption*{(a) TPC-DS Q11 Original}  % Sub-caption inside the same listing
  \end{minipage}
  \hfill
  \begin{minipage}[t]{0.49\linewidth}  % Second half
    \begin{tcolorbox}[colback=gray!5, colframe=gray!40, boxsep=2mm, left=6pt, right=5pt, top=0.5mm, bottom=0.5mm]
      \begin{minted}[
          frame=none,
          framesep=2mm,
          linenos,
          numbersep=2mm,
          escapeinside=||,
          baselinestretch=1.0, 
          fontsize=\tiny
      ]{sql}
with store_sales_total as (
        select ...
        from |\textcolor{blue}{$store\_sales$}|, ...
        ...
    ),
    web_sales_total as (
        select ...
        from |\textcolor{blue}{$web\_sales$}|, ...
        ...
    )
select
    s_2.customer_id, ...
from
    store_sales_total s_1,
    store_sales_total s_2,
    web_sales_total w_1,
    web_sales_total w_2
where
    s_2.customer_id = s_1.customer_id
    and s_1.customer_id = w_2.customer_id
    and s_1.customer_id = w_1.customer_id
    and s_1.dyear = 1998
    and w_1.dyear = 1998
    and s_2.dyear = 1999
    and w_2.dyear = 1999
    and w_2.y_total / w_1.y_total > 
        s_2.y_total / s_1.y_total
      \end{minted}
    \end{tcolorbox}
    \vspace{-1mm}
    \caption*{(b) TPC-DS Q11 Rewritten}
  \end{minipage}
  \vspace{-1mm}
  \captionsetup{font=small,name=Listing}
  \caption{Left: Original TPC-DS Q11 query. Right: A rewrite of Q11 generated by LLMs. The rewrite is 9x faster, though no existing rewrite rules can achieve this transformation.}
  \label{gen:listing:q11}
  \inv
\end{listing}
% \inv
\end{comment}



\begin{listing}[ht]
  \centering
  % --- Left Column: Original ---
  \begin{minipage}[t]{0.48\linewidth}
    \begin{framed}
      \vspace{-3mm} % Adjust spacing for framed environment
      \begin{lstlisting}
with year_total as (
    select ..., 's' sale_type
    from  |\textcolor{blue}{$store\_sales$}|, ...
    ...
    union all
    select ..., 'w' sale_type
    from  |\textcolor{blue}{$web\_sales$}|, ...
    ...
)
select
    s_2.customer_id, ...
from
    year_total s_1,
    year_total s_2,
    year_total w_1,
    year_total w_2
where
    s_2.customer_id = s_1.customer_id
    and s_1.customer_id = w_2.customer_id
    and s_1.customer_id = w_1.customer_id
    and s_1.sale_type = 's'
    and s_1.dyear = 1998
    and w_1.sale_type = 'w'
    and w_1.dyear = 1998
    and s_2.sale_type = 's'
    and s_2.dyear = 1999
    and w_2.sale_type = 'w'
    and w_2.dyear = 1999
    and w_2.y_total / w_1.y_total > 
        s_2.y_total / s_1.y_total
      \end{lstlisting}
      \vspace{-3mm}
    \end{framed}
    \vspace{-2mm}
    \caption*{(a) TPC-DS Q11 Original}
  \end{minipage}
  \hfill
  % --- Right Column: Rewritten ---
  \begin{minipage}[t]{0.48\linewidth}
    \begin{framed}
      \vspace{-3mm}
      \begin{lstlisting}
with store_sales_total as (
        select ...
        from |\textcolor{blue}{$store\_sales$}|, ...
        ...
    ),
    web_sales_total as (
        select ...
        from |\textcolor{blue}{$web\_sales$}|, ...
        ...
    )
select
    s_2.customer_id, ...
from
    store_sales_total s_1,
    store_sales_total s_2,
    web_sales_total w_1,
    web_sales_total w_2
where
    s_2.customer_id = s_1.customer_id
    and s_1.customer_id = w_2.customer_id
    and s_1.customer_id = w_1.customer_id
    and s_1.dyear = 1998
    and w_1.dyear = 1998
    and s_2.dyear = 1999
    and w_2.dyear = 1999
    and w_2.y_total / w_1.y_total > 
        s_2.y_total / s_1.y_total
      \end{lstlisting}
      \vspace{-3mm}
    \end{framed}
    \vspace{-2mm}
    \caption*{(b) TPC-DS Q11 Rewritten}
  \end{minipage}

  \vspace{1mm}
  \captionsetup{font=small}
  \caption{Left: Original TPC-DS Q11 query. Right: A rewrite of Q11 generated by LLMs. The rewrite is 9x faster, though no existing rewrite rules can achieve this transformation.}
  \label{gen:listing:q11}
\end{listing}





\genhead{Example 2: TPC-DS Q11} In Listing~\ref{gen:listing:q11} (left), a query adapted from TPC-DS Q11 identifies customers whose web sales growth from 1998 to 1999 outpaces their store sales growth over the same period. It begins by creating a Common Table Expression (CTE) \texttt{year\_total}, which aggregates annual store and web sales data for each customer (lines 2–4 and 6–8) via a \sqlunionallcolor operation (line 5). The query then filters this aggregated data by sales type and year to isolate the figures for 1998 and 1999 (lines 21–28). A join on \texttt{customer\_id} (lines 18–20) allows for a comparison of each customer's sales across different years and channels. In essence, this query uses \sqlunionallcolor to combine multiple data sources with a column to label each source, and then filtering on the label later to separate them. The combining and filtering cancel each other's effect and thus can be removed for better performance and readability.
 The revised query in Listing~\ref{gen:listing:q11} (Right) makes two changes: (1) creates separate CTEs for store and web sales data, and (2) removes the sale type filters. This results in a 9x speedup, reducing execution time from 18 minutes to 2 minutes on PostgreSQL with TPC-DS scale factor 10. The performance boost is due not only to avoiding unnecessary computations but also to PostgreSQL's cardinality estimation errors with more predicates.

No existing rewrite rules address this anti-pattern, preventing rule-based query rewriting techniques from discovering this more efficient rewrite. However, a database expert reviewing the example would intuitively reason that the \sqlunionallcolor operation and the subsequent filtering cancel each other out. However, manually scrutinizing every query for new anti-patterns is not feasible. LLMs with advanced reasoning capabilities offer a promising alternative for automatically identifying and exploiting new rewrite opportunities.


\begin{comment}
    


\begin{listing}[!ht]
\begin{tcolorbox}[colback=gray!5, colframe=gray!40, arc=0mm, outer arc=0mm, boxsep=2mm, left=15pt, right=5pt, top=0.5mm, bottom=0.5mm]
\begin{minted}[
    frame=none,
    framesep=2mm,
    baselinestretch=1.05,
    fontsize=\footnotesize,
    linenos,
    escapeinside=||
]{sql}
with year_total as (
    select ..., 's' sale_type
    from  |\textcolor{blue}{$store\_sales$}|, customer, date_dim
    ...
    union all
    select ..., 'w' sale_type
    from  |\textcolor{blue}{$web\_sales$}|, customer, date_dim
    ...
)
select
    s_99.customer_id,
    s_99.customer_email_address
from
    year_total s_98,
    year_total s_99,
    year_total w_98,
    year_total w_99
where
    s_99.customer_id = s_98.customer_id
    and s_98.customer_id = w_99.customer_id
    and s_98.customer_id = w_98.customer_id
    and s_98.sale_type = 's'
    and s_98.dyear = 1998
    and w_98.sale_type = 'w'
    and w_98.dyear = 1998
    and s_99.sale_type = 's'
    and s_99.dyear = 1998 + 1
    and w_99.sale_type = 'w'
    and w_99.dyear = 1998 + 1
    and w_99.y_total / w_98.y_total > 
        s_99.y_total / s_98.y_total
\end{minted}
\end{tcolorbox}
\vspace{-2mm}
\caption{TPC-DS Q11 aims to find the customers whose web sales growth rate from 1998 to 1999 is higher than their store sales growth rate over the same period.}
\label{gen:listing:1}
\end{listing}

\begin{listing}[ht]
\begin{tcolorbox}[colback=gray!5, colframe=gray!40, arc=0mm, outer arc=0mm, boxsep=2mm, left=15pt, right=5pt, top=0.5mm, bottom=0.5mm]
\begin{minted}[
    frame=none,
    framesep=2mm,
    baselinestretch=1.05,
    fontsize=\footnotesize,
    linenos,
    escapeinside=||
]{sql}
with store_sales_total as (
        select ...
        from |\textcolor{blue}{$store\_sales$}|, ...
        ...
    ),
    web_sales_total as (
        select ...
        from |\textcolor{blue}{$web\_sales$}|, ...
        ...
    )
select
    s_99.customer_id,
    s_99.customer_email_address
from
    store_sales_total s_98,
    store_sales_total s_99,
    web_sales_total w_98,
    web_sales_total w_99
where
    s_99.customer_id = s_98.customer_id
    and s_98.customer_id = w_99.customer_id
    and s_98.customer_id = w_98.customer_id
    and s_98.dyear = 1998
    and w_98.dyear = 1998
    and s_99.dyear = 1998 + 1
    and w_99.dyear = 1998 + 1
    and w_99.y_total / w_98.y_total > 
        s_99.y_total / s_98.y_total
\end{minted}
\end{tcolorbox}
\vspace{-2mm}
\caption{Rewrite TPC-DS Q11 to an equivalent query where two CTEs are created for store and web sales respectively.}
\label{gen:listing:2}
\end{listing}

\end{comment}






% \begin{listing}[!ht]
% \begin{tcolorbox}[colback=gray!5, colframe=gray!40, arc=0mm, outer arc=0mm, boxsep=2mm, left=15pt, right=5pt, top=0.5mm, bottom=0.5mm]
% \begin{lstlisting}[language=SQL, caption={TPC-DS Q11 aims to find the customers whose web sales growth rate from 1998 to 1999 is higher than their store sales growth rate over the same period.}, label=listing:1]
% with year_total as (
%     select ..., 's' sale_type
%     from  (*@\textcolor{blue}{\$store\_sales\$}@*), customer, date_dim
%     ...
%     union all
%     select ..., 'w' sale_type
%     from  (*@\textcolor{blue}{\$web\_sales\$}@*), customer, date_dim
%     ...
% )
% select
%     s_99.customer_id,
%     s_99.customer_email_address
% from
%     year_total s_98,
%     year_total s_99,
%     year_total w_98,
%     year_total w_99
% where
%     s_99.customer_id = s_98.customer_id
%     and s_98.customer_id = w_99.customer_id
%     and s_98.customer_id = w_98.customer_id
%     and s_98.sale_type = 's'
%     and s_98.dyear = 1998
%     and w_98.sale_type = 'w'
%     and w_98.dyear = 1998
%     and s_99.sale_type = 's'
%     and s_99.dyear = 1998 + 1
%     and w_99.sale_type = 'w'
%     and w_99.dyear = 1998 + 1
%     and w_99.y_total / w_98.y_total > 
%         s_99.y_total / s_98.y_total
% \end{lstlisting}
% \end{tcolorbox}
% \end{listing}


% \begin{listing}[ht]
% \begin{tcolorbox}[colback=gray!5, colframe=gray!40, arc=0mm, outer arc=0mm, boxsep=2mm, left=15pt, right=5pt, top=0.5mm, bottom=0.5mm]
% \begin{lstlisting}[caption={Rewrite TPC-DS Q11 to an equivalent query where two CTEs are created for store and web sales respectively.}, label=listing:2]
% with store_sales_total as (
%         select ...
%         from |\textcolor{blue}{\$store\_sales\$}|, ...
%         ...
%     ),
%     web_sales_total as (
%         select ...
%         from |\textcolor{blue}{\$web\_sales\$}|, ...
%         ...
%     )
% select
%     s_99.customer_id,
%     s_99.customer_email_address
% from
%     store_sales_total s_98,
%     store_sales_total s_99,
%     web_sales_total w_98,
%     web_sales_total w_99
% where
%     s_99.customer_id = s_98.customer_id
%     and s_98.customer_id = w_99.customer_id
%     and s_98.customer_id = w_98.customer_id
%     and s_98.dyear = 1998
%     and w_98.dyear = 1998
%     and s_99.dyear = 1998 + 1
%     and w_99.dyear = 1998 + 1
%     and w_99.y_total / w_98.y_total > 
%         s_99.y_total / s_98.y_total
% \end{lstlisting}
% \end{tcolorbox}
% \end{listing}


% and s_98.y_total > 0
% and w_98.y_total > 0






\begin{comment}

\begin{listing}[!ht]
\begin{tcolorbox}[colback=gray!5, colframe=gray!40, arc=0mm, outer arc=0mm, boxsep=2mm, left=15pt, right=5pt, top=0.5mm, bottom=0.5mm]
\begin{minted}[
    frame=none,
    framesep=2mm,
    baselinestretch=1.05,
    fontsize=\footnotesize,
    linenos,
    escapeinside=||
]{sql}
with year_total as (
    select
        c_customer_id customer_id,
        c_email_address customer_email_address,
        d_year dyear,
        sum(ss_ext_list_price - ss_ext_discount_amt) y_total,
        's' sale_type
    from  |\textcolor{blue}{$store\_sales$}|, customer, date_dim
    where
        c_customer_sk = ss_customer_sk
        and ss_sold_date_sk = d_date_sk
    group by  c_customer_id, c_email_address, d_year
    union all
    select
        c_customer_id customer_id,
        c_email_address customer_email_address,
        d_year dyear,
        sum(ws_ext_list_price - ws_ext_discount_amt) y_total,
        'w' sale_type
    from  |\textcolor{blue}{$web\_sales$}|, customer, date_dim
    where
        c_customer_sk = ws_bill_customer_sk
        and ws_sold_date_sk = d_date_sk
    group by  c_customer_id, c_email_address, d_year
)
select
    s_99.customer_id,
    s_99.customer_email_address
from
    year_total s_98,
    year_total s_99,
    year_total w_98,
    year_total w_99
where
    s_99.customer_id = s_98.customer_id
    and s_98.customer_id = w_99.customer_id
    and s_98.customer_id = w_98.customer_id
    and s_98.sale_type = 's'
    and s_98.dyear = 1998
    and w_98.sale_type = 'w'
    and w_98.dyear = 1998
    and s_99.sale_type = 's'
    and s_99.dyear = 1998 + 1
    and w_99.sale_type = 'w'
    and w_99.dyear = 1998 + 1
    and w_99.y_total / w_98.y_total > 
        s_99.y_total / s_98.y_total
\end{minted}
\end{tcolorbox}
\vspace{-2mm}
\caption{TPC-DS Q11 aims to find the customers whose web sales growth rate from 1998 to 1999 is higher than their store sales growth rate over the same period.}
\label{gen:listing:1}
\end{listing}










\begin{listing}[ht]
\begin{tcolorbox}[colback=gray!5, colframe=gray!40, arc=0mm, outer arc=0mm, boxsep=2mm, left=15pt, right=5pt, top=0.5mm, bottom=0.5mm]
\begin{minted}[
    frame=none,
    framesep=2mm,
    baselinestretch=1.05,
    fontsize=\footnotesize,
    linenos,
    escapeinside=||
]{sql}
with store_sales_total as (
        select ...
        from |\textcolor{blue}{$store\_sales$}|, ...
        ...
    ),
    web_sales_total as (
        select ...
        from |\textcolor{blue}{$web\_sales$}|, ...
        ...
    )
select
    s_99.customer_id,
    s_99.customer_email_address
from
    store_sales_total s_98,
    store_sales_total s_99,
    web_sales_total w_98,
    web_sales_total w_99
where
    s_99.customer_id = s_98.customer_id
    and s_98.customer_id = w_99.customer_id
    and s_98.customer_id = w_98.customer_id
    and s_98.dyear = 1998
    and w_98.dyear = 1998
    and s_99.dyear = 1998 + 1
    and w_99.dyear = 1998 + 1
    and w_99.y_total / w_98.y_total > 
        s_99.y_total / s_98.y_total
\end{minted}
\end{tcolorbox}
\vspace{-2mm}
\caption{Rewrite TPC-DS Q11 to an equivalent query where two CTEs are created for store and web sales respectively.}
\label{gen:listing:2}
\end{listing}


    
\end{comment}



\begin{comment}
    

The example query in Listing~\ref{gen:listing:1} finds customers with a higher web sales growth rate from 1998 to 1999 compared to their store sales growth rate. It first creates a Common Table Expression (CTE) \texttt{year\_total} combining store and web sales data for each customer and year (lines 2-8) using \sqlunionallcolor (line 5). Filters are then applied to extract data for the years 1998 and 1999 (lines 22-29), and a join is performed on customer id to compare sales (lines 19-21). Finally, it filters for customers with higher web sales growth (lines 30-31).
%%
 In essence, this query uses \sqlunionallcolor to combine multiple data sources with a column to label each source, and then filtering on the label later to separate them. The combining and filtering cancel each other's effect and thus can be removed for better performance and readability.
%% 
 The revised query in Listing~\ref{gen:listing:2} makes two changes: (1) creates separate CTEs for store and web sales data, and (2) removes the sale type filters. This results in a 9x speedup, reducing execution time from 18 minutes to 2 minutes on PostgreSQL with TPC-DS scale factor 10. The performance boost is due not only to avoiding unnecessary computations but also to PostgreSQL's cardinality estimation errors with more predicates. Additionally, the rewrite improves modularity and maintainability.

Creating a pattern-matching rule for this type of redundancy requires tracking label columns, and determining whether the subquery result is used elsewhere. It also requires rules for different clauses (e.g., WHERE, CASE WHEN). Given the variety of possible patterns, creating rules to eliminate all cases of  unnecessary \sqlunionallcolor will be daunting.
In contrast, a human expert would intuitively recognize that ``combining and then splitting'' is redundant. 
Manually scrutinizing every query is not scalable. 
However, 
LLMs, with their vast general knowledge, 
offer a promising alternative for
automating this process.

\end{comment}









Background on large language models, prompting, and token costs
appears in Section~\ref{sec:bg:tools}.
% NOTE (dissertation): the paper's LLM background subsection below is
% superseded by Chapter 2 and therefore skipped via \iffalse.
\iffalse
\subsection{Large Language Models: Background}

Recent advancements in deep learning and the availability of large datasets have propelled LLMs into the spotlight. These models  exhibit advanced reasoning skills, allowing them to take on open-ended tasks and make logical inferences. 

\genhead{Prompt engineering}
The key to leveraging LLMs effectively lies in \emph{prompt engineering}\cite{googleprompt}. This involves crafting precise inputs to guide the model's output. By providing natural language instructions, users can guide the model towards the desired outcome. Through iterative refinement, prompt engineering allows for high-quality and relevant responses, giving better control over creativity and consistency. In query rewriting, we concatenate an input query $q$ with a prompt $p$, denoted as $p \texttt{||} q$, which encodes a comprehensive task description for rewriting. Figure~\ref{gen:fig:basic_prompt} is an example of a basic prompt with only the query and a concise prompt,  which we refer to as \baseline.




\begin{comment}
\genhead{Zero-shot prompting vs few-shot prompting} The key to successfully using LLMs is to find the optimal prompt, commonly known as \emph{prompt engineering}\cite{googleprompt}. 
Prompt engineering is classified into two scenarios, depending on 
the number of examples provided in the prompt: zero-shot prompting\cite{gao2023text} and few-shot prompting~\cite{gu2023few}. In zero-shot prompting, no example is provided. Figure~\ref{gen:fig:basic_prompt} shows a simple prompt which only includes the input query itself and a concise one-sentence rewriting instruction. Conversely, in few-shot prompting, a limited number of examples are provided to the LLM, allowing it to identify explicit or implicit patterns from the input examples and generate corresponding output. For instance, in the context of Text-to-SQL, an example in the few-shot prompt may include a database schema, a natural language query and the corresponding SQL queries as a solution.

\end{comment}

\genhead{Tokens and cost} Tokens are the units or building blocks that LLMs use to process language. For example, the string `` tokenization'' is decomposed as two tokens `` token'' and ``ization'', while a short common word like `` the'' is represented as a single token~\cite{openaitokens}. The cost of using LLMs is typically determined by the number of tokens in both the input and the output.

% \vspace{2mm}
\inv






\fi
\begin{comment}


\subsection{Opportunities and Challenges of Introducing LLMs to Query Rewriting}

\genhead{Opportunities} Optimizing increasingly complex queries necessitates a comprehensive understanding of their semantics. Database experts have encoded their accumulated understanding and insights into an extensive set of rewrite rules over the past few decades. However, rules can never be sufficient as there will be always new scenarios and new query patterns that have not been encountered or envisioned. In instances where a rewrite opportunity eludes existing rules and human inspection is not feasible, LLMs can serve as a viable substitute for human experts. Specifically, LLMs excel in parsing lengthy and intricate SQL queries, extracting relevant general knowledge, and generating high-quality suggestions tailored to the query. In view of the recent advent and success of LLMs in performing complex tasks, leveraging their potential for query rewriting could lead to a revolution in identifying and exploiting new rewriting opportunities, particularly addressing various computational redundancies. This could significantly ease the burden on human experts who would otherwise need to manually inspect numerous slow queries, thereby enhancing overall productivity.

% \subsubsection{\textbf{Challenges}}
\genhead{Challenges} Despite their potential, effectively leveraging LLMs in query rewriting is challenging. (1) LLM responses are often untrustworthy, leading to a ``hallucination'' problem, while the correctness of the database remains paramount. It is imperative to design strategies that guide LLM to generate equivalent queries through reasonable analysis. Identifying errors and automating their correction with minimal human intervention becomes crucial. (2)
Depending on the query complexity and the frequency of relevant queries appearing in the training corpus, the effectiveness of using LLMs for rewriting may vary significantly. Some queries can be solved straightforwardly, while others require more guidance. There is a need to automatically find the best guidance for such queries. Developing an effective method to extract knowledge from successful instances for further reuse is essential to maximize the success rate.

% Due to the inherent randomness of the model, LLMs may be able to find significantly faster rewrites for some queries but fall short in optimizing others. Developing an effective method to extract knowledge from the rewriting processes for further reuse is essential to maximize the success rate.

\end{comment}


\begin{comment}

Despite the abundance of rewrite rules crafted by human experts over decades, many queries still fail to be rewritten into a more efficient form by the existing query optimizers. The objective of query optimization is to generate good execution plans for complex queries within a relatively short optimization time. As the number of rewrite rules increases, the search space expands rapidly, leading to a corresponding increase in optimization time. Query optimizers opt to forego certain rewrite opportunities in favor of short optimization time. Adding to the complexity is the fact that the applying rewrite rules does not always result in speedup; some rewrites may prove ineffective or even lead to a slowdown in query performance. For example, group-by pushdown transformation pushes the group-by operator down past joins in the query, thereby performing an early aggregation. Doing an early aggregation may result in a significant reduction in the number of rows later used in the join, hence the overall performance of the query may improve. On the other hand, if the joins in the original query significantly reduce the size of the data to be aggregated, this rewrite could result in a slower execution plan. These trade-offs further complicate the optimization process. Consequently, query optimizers tend to consider only a small subset of the rules with high possibility to improve the query performance.

Simply adding more rewrite rules does not effectively address all issues, as it is inherently challenging for pattern-matching rules to capture certain rewrite opportunities. An optimization strategy that is intuitive to humans may correspond to numerous possible patterns, posing a difficulty in comprehensive coverage through pattern-matching rewrite rules. 

\end{comment}





\begin{comment}

\begin{listing}[!ht]
\begin{tcolorbox}[colback=gray!5, colframe=gray!40, arc=0mm, outer arc=0mm, boxsep=2mm, left=15pt, right=5pt, top=0.5mm, bottom=0.5mm]
\begin{minted}[
    frame=none,
    framesep=2mm,
    baselinestretch=1.05,
    fontsize=\tiny,
    linenos,
    escapeinside=||
]{sql}
with year_total as (
    select ..., 's' sale_type
    from  |\textcolor{blue}{$store\_sales$}|, customer, date_dim
    ...
    union all
    select ..., 'w' sale_type
    from  |\textcolor{blue}{$web\_sales$}|, customer, date_dim
    ...
)
select
    s_99.customer_id,
    s_99.customer_email_address
from
    year_total s_98,
    year_total s_99,
    year_total w_98,
    year_total w_99
where
    s_99.customer_id = s_98.customer_id
    and s_98.customer_id = w_99.customer_id
    and s_98.customer_id = w_98.customer_id
    and s_98.sale_type = 's'
    and s_98.dyear = 1998
    and w_98.sale_type = 'w'
    and w_98.dyear = 1998
    and s_99.sale_type = 's'
    and s_99.dyear = 1998 + 1
    and w_99.sale_type = 'w'
    and w_99.dyear = 1998 + 1
    and w_99.y_total / w_98.y_total > 
        s_99.y_total / s_98.y_total
\end{minted}
\end{tcolorbox}
\vspace{-2mm}
\caption{TPC-DS Q11 aims to find the customers whose web sales growth rate from 1998 to 1999 is higher than their store sales growth rate over the same period.}
\label{gen:listing:1}
\end{listing}

\begin{listing}[ht]
\begin{tcolorbox}[colback=gray!5, colframe=gray!40, arc=0mm, outer arc=0mm, boxsep=2mm, left=13pt, right=3pt, top=0.5mm, bottom=0.5mm]
\begin{minted}[
    frame=none,
    framesep=1mm,
    baselinestretch=1.05,
    fontsize=\tiny,
    linenos,
    escapeinside=||
]{sql}
with store_sales_total as (
        select ...
        from |\textcolor{blue}{$store\_sales$}|, ...
        ...
    ),
    web_sales_total as (
        select ...
        from |\textcolor{blue}{$web\_sales$}|, ...
        ...
    )
select
    s_99.customer_id,
    s_99.customer_email_address
from
    store_sales_total s_98,
    store_sales_total s_99,
    web_sales_total w_98,
    web_sales_total w_99
where
    s_99.customer_id = s_98.customer_id
    and s_98.customer_id = w_99.customer_id
    and s_98.customer_id = w_98.customer_id
    and s_98.dyear = 1998
    and w_98.dyear = 1998
    and s_99.dyear = 1998 + 1
    and w_99.dyear = 1998 + 1
    and w_99.y_total / w_98.y_total > 
        s_99.y_total / s_98.y_total
\end{minted}
\end{tcolorbox}
\vspace{-2mm}
\caption{Rewrite TPC-DS Q11 to an equivalent query where two CTEs are created for store and web sales respectively.}
\label{gen:listing:2}
\end{listing}


\end{comment}


\begin{comment}

In fact, the \sqlunionallcolor and the filtering operations cancel each other's effect and should be removed from the query for better performance and readability. The rewrite, as presented in Listing~\ref{gen:listing:q11} (right), only makes two changes: (1) generates two CTEs for store and web sales data separately instead of combining them into one using \sqlunionallcolor, (2) removes filters regarding sale type in the main query. Consequently, the observed speedup, when tested on PostgreSQL with TPC-DS scale factor set to 10, is approximately 10.3x, reducing the execution time from 20 minutes to 2 minutes. The significant speedup is not solely attributed to avoiding all computations associated with \sqlunionallcolor and predicates on the sale type column. Rather, the observed performance gap is mainly attributed to the fact that in PostgreSQL the trend to underestimate cardinality is more pronounced with more predicates in the query. Clearly, the original query has more predicates than the rewrite and PostgreSQL chose nested-loop joins over sort-merge joins for the original query due to the severe cardinality underestimates. In addition to the performance boost, this rewrite also enhances modularity and maintainability.

The query on the left in Listing~\ref{gen:listing:q11}, derived from Q11 of the TPC-DS benchmark, aims to find the customers whose web sales growth rate from 1998 to 1999 is higher than their store sales growth rate over the same period. The query first creates a Common Table Expression (CTE) \texttt{year\_total} that combines the store and web sales data for each customer for each year (lines 2-4 and lines 6-8) using \sqlunionallcolor  (line 5). It then applies filters to the sale type and year columns to extract store and web sales data for each customer specifically for the years 1998 and 1999 (lines 22-29) and joins the results on customer id to facilitates a comparison of sales data for the same customer across different years and sale types (lines 19-21). Finally, it filters for customers whose web sales growth rate is higher than their store sales growth rate (lines 30-31). More generally, a poorly written query may employ \sqlunionallcolor to combine multiple data sources, using a single column to label each part, and then apply filters on the label column to use each part separately. In such cases, the combining and the filtering operations cancel each other's effect and should be removed from the query for better performance and readability. The revised query, as presented in Listing~\ref{gen:listing:2}, only makes two changes: (1) generates two CTEs for store and web sales data separately instead of combining them into one, (2) removes filters regarding sale type in the main query. Consequently, the observed speedup, when tested on PostgreSQL with TPC-DS scale factor set to 10, is approximately 9x, reducing the execution time from 18 minutes to 2 minutes. The significant speedup is not solely attributed to avoiding all computations associated with \sqlunionallcolor and predicates on the sale type column. Rather, the observed performance gap is also attributed to the fact that in PostgreSQL the trend to underestimate cardinalities is more pronounced with more predicates in the query. Clearly, the original query has more predicates than the revised one and PostgreSQL chose nested-loop joins over sort-merge joins for the original query due to the severe cardinality underestimates. In addition to the performance boost, this rewrite also enhances modularity and maintainability.

    
\end{comment}


\begin{comment}


To address this kind of computational redundancy through pattern-matching rules, the initial challenge lies in identifying the label column(s). Then it is essential to track how the subquery's result is utilized and determine whether the combined result remains unused. This usage may occur in predicate clauses, case when clauses, or essentially any part within the outer query blocks. Given the multitude of possible patterns, we may still fall short of covering all cases for eliminating the unnecessary \sqlunionallcolor despite devising numerous rules. However, a human observer, when presented with the above example, would intuitively recognize that the process of ``combining and then splitting'' actually does nothing and should be removed from the query. This intuition stems from general knowledge rather than explicit rules. Given the limited time and resources of human experts, manually scrutinizing every query for potential redundancies is not feasible. While human wisdom remains invaluable, large language models (LLMs) like GPT-4 and beyond, which encode vast amounts of general knowledge, could be a potential game-changer, enabling the automatic identification and exploitation of rewrite opportunities.


\end{comment}


\begin{comment}
    This section provides a brief overview of LLM-related terminology used in the rest of this paper.

\genhead{Tokens} Tokens are the units or building blocks that LLMs use to process language. For example, the string `` tokenization'' is decomposed as two tokens `` token'' and ``ization'', while a short common word like `` the'' is represented as a single token~\cite{openaitokens}. The cost of using LLMs is typically determined by the number of tokens in both the input and the output.



\end{comment}